\documentclass[11pt]{article}

\usepackage[T1]{fontenc}
\usepackage[utf8]{inputenc}
\usepackage[a4paper,margin=1in]{geometry}
\usepackage{booktabs}
\usepackage{graphicx}
\usepackage{xcolor}
\usepackage{titlesec}
\usepackage{enumitem}
\usepackage[hidelinks]{hyperref}

\definecolor{accent}{HTML}{C0392B}
\definecolor{blu}{HTML}{2C6FBB}
\definecolor{ink}{HTML}{2B2B2B}

\titleformat{\section}{\Large\bfseries\color{ink}}{\thesection}{0.6em}{}
\titleformat{\subsection}{\large\bfseries\color{blu}}{\thesubsection}{0.6em}{}
\setlist{itemsep=2pt,topsep=3pt}
\graphicspath{{figures/}}

\newcommand{\artist}[1]{\textit{#1}}

\begin{document}

\begin{titlepage}
\centering
\vspace*{3cm}
{\color{accent}\rule{\textwidth}{2pt}}\\[0.8em]
{\Huge\bfseries A Portrait of a Listener}\\[0.4em]
{\Large Your music taste, read from the machine}\\[0.6em]
{\large\itshape Un retrato de quien escucha --- tu gusto musical, leído por la máquina}\\[0.8em]
{\color{accent}\rule{\textwidth}{2pt}}\\[2cm]
{\large A data portrait derived from 23{,}529 tracks of listening history,\\
the concepts a neural taste-model learned, and GPT-labeled sparse-autoencoder atoms.}\\[1cm]
{\normalsize Spotify Engagement Lab \quad$\cdot$\quad rotation dataset \texttt{ds-20260612-143014-p64-s42}}\\[0.3em]
{\normalsize 2026-07-12}
\vfill
\end{titlepage}

% ============================================================
\section*{How to read this document}
This report is bilingual. \textbf{Part~I is in English}; \textbf{Part~II es la
versión en español} (identical content). It is built from three sources: (1)~the
raw listening history of a single Spotify user (23{,}529 tracks); (2)~two neural
networks trained to predict which tracks that user \emph{replays}
(\texttt{rotation} $=$ played $\geq 2$ times), the best taste-models on record;
and (3)~a sparse-autoencoder (SAE) probe of those models' hidden layers, whose
learned ``atoms'' were labeled by GPT (\texttt{gpt-4o-mini}). Method and caveats
are in the Provenance appendix --- read it: the model captures a \emph{real but
modest} slice of taste, and the honest picture is as interesting as the flattering
one.

\clearpage
\begin{center}{\color{blu}\Large\bfseries PART I \quad$\cdot$\quad ENGLISH}\end{center}
\vspace{0.5em}

\section{The one-paragraph portrait}
You are, above all, an \textbf{electronic listener with a rock heart and an
Argentine passport}. Of the 23{,}529 tracks in your history you come back to
\textbf{47.3\%} of them --- a high replay rate for a library this eclectic. What
you \emph{return} to skews electronic and produced: intelligent dance, downtempo,
IDM, deep and organic house. But the artists you play most are as likely to be
guitar bands --- German industrial, British art-rock, and a deep seam of Argentine
rock and indie. You favour the \textbf{instrumental over the lyrical}, the
\textbf{moody over the sunny}, and \textbf{curated mid-fame artists over chart
stars}. You are not a genre loyalist; you are a texture loyalist.

\section{What you actually replay}
\subsection{The artists you come back to}
Ranked by total plays among replayed tracks:

\begin{center}
\begin{tabular}{@{}rl@{\hskip 2em}rl@{}}
\toprule
\# & Artist & \# & Artist \\
\midrule
1 & \artist{Rammstein} & 11 & \artist{Caribou} \\
2 & \artist{Babasónicos} & 12 & \artist{Massive Attack} \\
3 & \artist{Peces Raros} & 13 & \artist{Laika Perra Rusa} \\
4 & \artist{Kraftwerk} & 14 & \artist{Daft Punk} \\
5 & \artist{Four Tet} & 15 & \artist{Tosca} \\
6 & \artist{Aphex Twin} & 16 & \artist{Air} \\
7 & \artist{Underworld} & 17 & \artist{Thievery Corporation} \\
8 & \artist{The Avalanches} & 18 & \artist{Radiohead} \\
9 & \artist{Usted Señálemelo} & 19 & \artist{Gorillaz} \\
10 & \artist{Joy Orbison} & 20 & \artist{Daphni} \\
\bottomrule
\end{tabular}
\end{center}

\subsection{Three pillars}
\begin{description}[leftmargin=1.4em,style=nextline]
\item[Electronic auteurs \& intelligent dance.] \artist{Four Tet}, \artist{Aphex
Twin}, \artist{Caribou}/\artist{Daphni}, \artist{Underworld}, \artist{Daft Punk},
\artist{Kraftwerk}, \artist{Joy Orbison}, \artist{Tosca}, \artist{Thievery
Corporation}, \artist{The Avalanches}. Downtempo, big beat, IDM, deep/organic
house, Düsseldorf electronic. This is the centre of gravity.
\item[Argentine rock \& indie.] \artist{Babasónicos}, \artist{Peces Raros},
\artist{Usted Señálemelo}, \artist{Laika Perra Rusa}. A distinct national thread ---
argentine indie, argentine rock, argentine alternative --- woven right through the
electronic core.
\item[Alt-rock crossover.] \artist{Radiohead}, \artist{Gorillaz},
\artist{Massive Attack}, \artist{Rammstein}. Rock that shakes hands with
electronics --- trip-hop, art-rock, industrial.
\end{description}

\subsection{Genre families: what earns a replay}
Not every genre earns its place in the library. Measuring \emph{replay lift} --- a
family's share of your \emph{replays} divided by its share of the whole library ---
electronic/dance is the clear winner, while hip-hop and jazz are more likely to be
heard once and set aside.

\begin{center}\includegraphics[width=0.86\textwidth]{taste_genre_lift.pdf}\end{center}

\section{Your acoustic fingerprint}
Comparing the average replayed track to the average play-once-and-drop track
isolates what, sonically, earns a second listen. The single strongest signal is
\textbf{instrumentalness}: the tracks you return to are markedly more instrumental
--- consistent with an electronic/downtempo core where the voice is texture, not
message. Replays are also a touch more energetic and danceable, and slightly
\emph{less} acoustic and \emph{less} bright (lower valence): produced, propulsive,
faintly melancholic.

\begin{center}\includegraphics[width=0.86\textwidth]{taste_acoustic.pdf}\end{center}

\section{When your music was made}
Your replayed music centres on the \textbf{2010s} (median release year
\textbf{2014}) but reaches back through the 2000s and 1990s and touches the
Kraftwerk-era 1970s. Modern-leaning, but not novelty-driven.

\begin{center}\includegraphics[width=0.78\textwidth]{taste_era.pdf}\end{center}

You also replay \textbf{mid-fame} artists (median artist popularity 55/100,
$\sim$390k followers) rather than chart stars --- a curatorial, slightly-off-centre
palette rather than a mainstream one.

\section{What the machine sees}
The neural taste-models don't store ``you like Four Tet.'' They learn abstract
\emph{feature detectors} in their hidden layers; a sparse autoencoder pulls those
apart into individual ``atoms,'' and GPT named each one from the tracks that most
excite it. Two honest truths emerged. First, \textbf{your taste is not cleanly
separable} --- most atoms are \emph{polysemantic}, blending genres, because what
unites your replays is texture and identity more than any one label. Second, the
concepts that best predict a \emph{replay} are exactly the portrait above.

\subsection{Concepts that predict a replay ($+$)}
\begin{center}
\begin{tabular}{@{}rll@{}}
\toprule
corr & Concept (GPT label) & Model \\
\midrule
$+0.39$ & Chillwave \& organic house & mlp \\
$+0.36$ & Bossa nova \& deep house & mlp \\
$+0.33$ & Upbeat electronic (reggaeton, deep house) & emb \\
$+0.32$ & Electronic classics & mlp \\
$+0.32$ & Mellow indie \& dance & emb \\
$+0.30$ & Argentine rock favourites & emb \\
$+0.29$ & Cumbia \& organic house & emb \\
$+0.26$ & Chill \& downtempo & emb \\
$+0.26$ & Argentine rock \& ambient & emb \\
\bottomrule
\end{tabular}
\end{center}

\subsection{Concepts that predict a skip ($-$)}
\begin{center}
\begin{tabular}{@{}rll@{}}
\toprule
corr & Concept (GPT label) & Model \\
\midrule
$-0.31$ & Folk \& alternative-rock (dropped) & mlp \\
$-0.24$ & Electro-swing \& indie (dropped) & mlp \\
$-0.24$ & Dance-pop \& indie (dropped) & mlp \\
$-0.16$ & ``Unpopular'' one-and-done tracks & emb \\
$-0.13$ & Indie \& chillwave you skip & emb \\
$-0.12$ & Dance-pop \& alt-metal you skip & emb \\
\bottomrule
\end{tabular}
\end{center}

The negatives are as telling as the positives: \textbf{dance-pop, electro-swing
and folk} are the sounds most likely to get one play and no more.

\section{The portrait in one line}
\begin{quote}\itshape
Instrumental, produced, slightly melancholic electronic music --- intelligent dance
and downtempo at the core, a spine of Argentine rock and indie, and an alt-rock
crossover streak --- replayed from mid-fame, curatorial artists, mostly from the
2010s. Dance-pop and folk need not apply.
\end{quote}

\clearpage
\begin{center}{\color{blu}\Large\bfseries PARTE II \quad$\cdot$\quad ESPAÑOL}\end{center}
\vspace{0.5em}

\section{El retrato en un párrafo}
Eres, ante todo, \textbf{alguien de la electrónica, con corazón de rock y
pasaporte argentino}. De los 23{,}529 temas de tu historial vuelves al
\textbf{47{,}3\,\%} --- una tasa de reescucha alta para una biblioteca tan
ecléctica. A lo que \emph{regresas} tiende a lo electrónico y producido:
\emph{intelligent dance}, \emph{downtempo}, IDM, \emph{deep} y \emph{organic
house}. Pero los artistas que más suenas son con la misma frecuencia bandas de
guitarras --- industrial alemán, art-rock británico y una veta profunda de rock e
indie argentinos. Prefieres lo \textbf{instrumental sobre lo lírico}, lo
\textbf{melancólico sobre lo luminoso}, y \textbf{artistas de fama media y
curados por sobre las estrellas de las listas}. No eres leal a un género; eres
leal a una textura.

\section{Lo que realmente reescuchas}
\subsection{Los artistas a los que vuelves}
Ordenados por reproducciones totales entre los temas reescuchados:

\begin{center}
\begin{tabular}{@{}rl@{\hskip 2em}rl@{}}
\toprule
\# & Artista & \# & Artista \\
\midrule
1 & \artist{Rammstein} & 11 & \artist{Caribou} \\
2 & \artist{Babasónicos} & 12 & \artist{Massive Attack} \\
3 & \artist{Peces Raros} & 13 & \artist{Laika Perra Rusa} \\
4 & \artist{Kraftwerk} & 14 & \artist{Daft Punk} \\
5 & \artist{Four Tet} & 15 & \artist{Tosca} \\
6 & \artist{Aphex Twin} & 16 & \artist{Air} \\
7 & \artist{Underworld} & 17 & \artist{Thievery Corporation} \\
8 & \artist{The Avalanches} & 18 & \artist{Radiohead} \\
9 & \artist{Usted Señálemelo} & 19 & \artist{Gorillaz} \\
10 & \artist{Joy Orbison} & 20 & \artist{Daphni} \\
\bottomrule
\end{tabular}
\end{center}

\subsection{Tres pilares}
\begin{description}[leftmargin=1.4em,style=nextline]
\item[Autores de la electrónica e \emph{intelligent dance}.] \artist{Four Tet},
\artist{Aphex Twin}, \artist{Caribou}/\artist{Daphni}, \artist{Underworld},
\artist{Daft Punk}, \artist{Kraftwerk}, \artist{Joy Orbison}, \artist{Tosca},
\artist{Thievery Corporation}, \artist{The Avalanches}. \emph{Downtempo},
\emph{big beat}, IDM, \emph{deep}/\emph{organic house}, electrónica de
Düsseldorf. Es tu centro de gravedad.
\item[Rock e indie argentinos.] \artist{Babasónicos}, \artist{Peces Raros},
\artist{Usted Señálemelo}, \artist{Laika Perra Rusa}. Una hebra nacional bien
marcada --- indie, rock y alternativo argentinos --- entretejida en el núcleo
electrónico.
\item[Cruce con el rock alternativo.] \artist{Radiohead}, \artist{Gorillaz},
\artist{Massive Attack}, \artist{Rammstein}. Rock que le da la mano a la
electrónica: \emph{trip-hop}, art-rock, industrial.
\end{description}

\subsection{Familias de género: qué se gana una reescucha}
No todos los géneros se ganan su lugar. Midiendo el \emph{empuje de reescucha} ---
la participación de una familia en tus \emph{reescuchas} dividida por su
participación en toda la biblioteca --- la electrónica/dance gana con claridad,
mientras que el hip-hop y el jazz tienden a escucharse una vez y dejarse de lado.

\begin{center}\includegraphics[width=0.86\textwidth]{taste_genre_lift.pdf}\end{center}
\begin{center}\small\itshape (Etiquetas de los ejes en inglés en la figura: ``replay lift''
$=$ empuje de reescucha.)\end{center}

\section{Tu huella acústica}
Comparar el tema reescuchado promedio con el que se toca una vez y se descarta
aísla qué, sonoramente, se gana una segunda escucha. La señal más fuerte es la
\textbf{instrumentalidad}: los temas a los que vuelves son marcadamente más
instrumentales --- coherente con un núcleo electrónico/\emph{downtempo} donde la
voz es textura, no mensaje. Las reescuchas también son algo más enérgicas y
bailables, y un poco \emph{menos} acústicas y \emph{menos} luminosas (menor
\emph{valence}): producidas, propulsivas, levemente melancólicas.

\begin{center}\includegraphics[width=0.86\textwidth]{taste_acoustic.pdf}\end{center}

\section{De qué época es tu música}
Tu música reescuchada se centra en los \textbf{años 2010} (año de lanzamiento
mediano \textbf{2014}) pero se extiende por los 2000 y los 90 y toca los 70 de la
era Kraftwerk. Con inclinación a lo moderno, pero no por novedad.

\begin{center}\includegraphics[width=0.78\textwidth]{taste_era.pdf}\end{center}

Además reescuchas artistas de \textbf{fama media} (popularidad mediana 55/100,
$\sim$390k seguidores) antes que estrellas de listas --- una paleta curada y algo
descentrada más que \emph{mainstream}.

\section{Lo que ve la máquina}
Los modelos neuronales no guardan ``te gusta Four Tet''. Aprenden
\emph{detectores de rasgos} abstractos en sus capas ocultas; un
\emph{autoencoder} disperso (SAE) los separa en ``átomos'' individuales, y GPT
nombró cada uno a partir de los temas que más lo activan. Surgieron dos verdades
honestas. Primera: tu gusto \textbf{no es claramente separable} --- la mayoría de
los átomos son \emph{polisemánticos}, mezclan géneros, porque lo que une tus
reescuchas es más la textura y la identidad que una sola etiqueta. Segunda: los
conceptos que mejor predicen una \emph{reescucha} son exactamente el retrato de
arriba.

\subsection{Conceptos que predicen una reescucha ($+$)}
\begin{center}
\begin{tabular}{@{}rll@{}}
\toprule
corr & Concepto (etiqueta GPT) & Modelo \\
\midrule
$+0.39$ & \emph{Chillwave} y \emph{organic house} & mlp \\
$+0.36$ & Bossa nova y \emph{deep house} & mlp \\
$+0.33$ & Electrónica animada (reggaetón, \emph{deep house}) & emb \\
$+0.32$ & Clásicos electrónicos & mlp \\
$+0.32$ & Indie suave y dance & emb \\
$+0.30$ & Favoritos del rock argentino & emb \\
$+0.29$ & Cumbia y \emph{organic house} & emb \\
$+0.26$ & \emph{Chill} y \emph{downtempo} & emb \\
$+0.26$ & Rock argentino y \emph{ambient} & emb \\
\bottomrule
\end{tabular}
\end{center}

\subsection{Conceptos que predicen un descarte ($-$)}
\begin{center}
\begin{tabular}{@{}rll@{}}
\toprule
corr & Concepto (etiqueta GPT) & Modelo \\
\midrule
$-0.31$ & Folk y rock alternativo (descartados) & mlp \\
$-0.24$ & \emph{Electro-swing} e indie (descartados) & mlp \\
$-0.24$ & \emph{Dance-pop} e indie (descartados) & mlp \\
$-0.16$ & Temas ``poco populares'' de una sola escucha & emb \\
$-0.13$ & Indie y \emph{chillwave} que salteas & emb \\
$-0.12$ & \emph{Dance-pop} y \emph{alt-metal} que salteas & emb \\
\bottomrule
\end{tabular}
\end{center}

Los negativos dicen tanto como los positivos: \textbf{dance-pop, electro-swing y
folk} son los sonidos con más probabilidad de recibir una sola reproducción.

\section{El retrato en una línea}
\begin{quote}\itshape
Música electrónica instrumental, producida y algo melancólica --- \emph{intelligent
dance} y \emph{downtempo} en el núcleo, una columna de rock e indie argentinos y
una veta de cruce con el rock alternativo --- reescuchada de artistas de fama media
y curados, sobre todo de los años 2010. El \emph{dance-pop} y el folk que no se
molesten en aplicar.
\end{quote}

\clearpage
\section*{Provenance \& honest caveats \quad$\cdot$\quad Procedencia y advertencias}

\textbf{Data.} A single user's Spotify Extended Streaming History, materialized as
the rotation dataset \texttt{ds-20260612-143014-p64-s42} (23{,}529 tracks; target
\texttt{rotation} $=$ play\_count $\geq 2$; 11{,}127 replayed / 47.3\%). Aggregate
statistics (artists, genres, acoustics, era) are computed \emph{directly} over
this data and are exact.

\textbf{Models.} Two \texttt{burn-deep-classifier} neural nets, the best
taste-models on record on this dataset:
\texttt{emb-128-64-ed8-highvocab} (topology embeddings, AUC \textbf{0.72274}) and
\texttt{mlp-256-128-64-highvocab} (topology mlp, AUC 0.71800). They predict a
binary replay signal from song metadata (artist/genre one-hots $+$ acoustic PCA).

\textbf{Interpretability.} Per-layer sparse autoencoders (2$\times$ overcomplete,
$\ell_1=0.0015$, 40 epochs) trained on each model's hidden activations; analyses
\texttt{interp-run-20260712-171243-f65f7-msae} (emb) and
\texttt{interp-run-20260712-171915-c52e4-msae} (mlp). The atoms surfaced here are
the ones most correlated with the replay target (12 per layer); their top-12
activating tracks were sent to \textbf{GPT (\texttt{gpt-4o-mini})} for a concept
label and description.

\textbf{The honest caveat.} This model captures a \emph{real but modest} slice of
your taste. On a held-out split the replay signal is only weakly decodable from
the hidden layers (linear $R^2 \approx 0.14$), and the SAE found \emph{no clean,
single-genre ``taste atom''} --- the atoms are polysemantic, and a dataset-level
probe surfaced zero monosemantic target atoms at this data scale. In plain terms:
your taste is driven by \emph{artist and texture identity} more than by any tidy
genre rule the model could isolate, so the model-derived concepts (\S6) are
suggestive, while the aggregate portrait (\S2--\S5) is the solid ground. The
GPT labels are automated readings of 12 tracks each and occasionally blur; they
are illustrative, not definitive.

\vfill
\begin{center}\small
Generated from the Spotify Engagement Lab
$\cdot$ campaign \texttt{2026-07-12-high-vocab-capacity-scan}
$\cdot$ figures: \texttt{docs/figures/make\_taste\_figures.py}
\end{center}

\end{document}
